% !TeX root = ../main.tex
% !Mode:: "TeX:UTF-8"
\section{第二章 项目简介}\label{ux7b2cux4e8cux7ae0-ux9879ux76eeux7b80ux4ecb}

\subsection{2.1 命题内容}\label{ux547dux9898ux5185ux5bb9}

人工智能正在从提供信息和辅助决策，向参与实际业务执行发展。代码修复、功能开发、数据分析等场景，不仅需要模型给出答案，还需要智能体调用工具、执行操作并完成任务。Gartner 于 2025 年发布预测，到 2026 年底，40\% 的企业应用将集成面向特定任务的智能体，而发布预测时这一比例不足 5\%。智能体应用范围的扩大，也将带动对配套运行环境和基础服务的需求。

然而，应用需求的增长并不意味着智能体能够直接实现规模化落地。Gartner 指出，部署成本上升、商业价值不明确以及风险控制不足，是智能体项目推进过程中面临的重要障碍。从企业实际使用角度看，智能体的价值不仅取决于单次任务能否完成，还取决于能否同时服务多个用户、持续稳定运行，以及投入与产出是否匹配。因此，支撑智能体运行的基础软件，需要在功能、安全、效率和成本之间形成适合企业使用的平衡。

与此同时，云计算为智能体应用提供了集中部署和规模化服务的基础。中国信息通信研究院《云计算蓝皮书（2025年）》列示，2024 年全球云计算市场规模达到 6929 亿美元，人工智能的发展正在持续带动云计算需求。 在这一产业背景下，提高已有基础设施的任务承载能力，成为智能体运行服务的重要发展方向。华为技术有限公司围绕国产操作系统软件生态，提出``高密部署智能体沙箱引擎''命题，要求在有限资源条件下，为更多智能体提供独立、可靠的任务执行环境。

本课题旨在探索适合智能体任务特点的新机制、新方法和基础软件，提升有限硬件条件下的智能体部署与服务能力。项目基于企业指定的 openEuler/Conch + StratoVirt 环境运行 OpenClaw，在管理面与全部容器总内存限制为 16GB、每个沙箱内部可见规格为 2U8G、关闭 Swap 的条件下，支持代码修复、功能开发等指定任务。在保障任务正确完成、不同任务互不干扰的基础上，提高能够同时运行的智能体数量，为企业以可控投入扩大智能体应用规模提供支撑。

\subsection{2.2 项目概述}\label{ux9879ux76eeux6982ux8ff0}

PI OS 是面向企业研发团队与智能体服务平台的高密部署智能体沙箱引擎。项目依托华为与天津大学的校企合作研究基础，围绕代码修复、功能开发等应用场景，为智能体提供独立的任务执行空间，支持多个智能体同时开展工作，保障不同任务之间互不干扰。项目已完成代码开发，正围绕企业规定的环境开展测试优化，致力于为企业提供稳定、可靠、便于部署的智能体运行基础。

PI OS 以``让同样的硬件完成更多任务''为核心目标，将资源利用效率与企业实际使用需求相结合。在保障任务正常执行的前提下，通过更合理地组织运行环境和分配资源，提高有限服务器的任务承载能力，减少企业对持续扩充硬件的依赖。项目不仅关注智能体``能否运行''，更关注其能否以可控成本持续服务更多用户，力求推动智能体从单个任务试用走向企业规模化应用，为国产操作系统生态下的智能体服务提供可行的产品方案。

\subsection{2.3 项目进度}\label{ux9879ux76eeux8fdbux5ea6}

目前，PI OS 已完成代码开发，并形成系统设计材料和阶段测试报告。前期研究环境下，已开展 48 个任务的并发测试，报告记录任务全部完成，未发生内存不足导致的中断；在相同 48 个任务的补充测试中，总完成时间由约 735 秒缩短至 245 秒，缩短约 66.7\%。下一阶段将围绕华为指定运行环境开展适配与验证，进一步确认在 16GB 总内存限制下的智能体承载能力，完善部署工具和交付文档，为企业试点与产品转化奠定基础。

\subsection{2.4 项目使命}\label{ux9879ux76eeux4f7fux547d}

PI OS 的使命，不仅是让有限的硬件承载更多智能体，更是立足国产基础软件生态，以自主创新服务国家科技自立自强。项目坚持从真实需求出发，让智能体用得起、用得好、用得安心，帮助企业降低应用门槛，为科研探索和产业创新提供可靠支撑。让有限的计算资源服务更多创新，让技术进步惠及更多用户，为我国数字经济发展贡献青年力量。

\subsection{2.5 项目业务简介}\label{ux9879ux76eeux4e1aux52a1ux7b80ux4ecb}

PI OS 的目标客户首先是具有智能体应用需求的企业研发团队和高校科研团队，优先围绕代码修复、功能开发等场景开展试点。在完成产品验证并明确商业化授权后，逐步面向智能体服务平台、云服务商及其他有需求的企业推广。

PI OS 为客户提供智能体运行与管理能力，支持多个任务同时开展，并为不同任务提供相互独立的执行空间。产品以提高有限硬件的任务承载能力为核心，帮助客户更充分地利用已有服务器，兼顾任务完成效率、运行稳定性和使用成本。

客户完成产品部署与任务接入后，即可利用 PI OS 支撑智能体执行工作，并查看任务结果与资源使用情况。项目拟配套提供部署指南、使用文档、技术支持和持续维护服务，降低客户的接入与维护负担，逐步形成从试点验证到正式交付、长期服务的业务模式。

\subsection{2.6 战略目标}\label{ux6218ux7565ux76eeux6807}

PI OS 力争在十年内成长为具有国际竞争力的国产智能体运行基础软件品牌，成为企业规模化部署智能体的重要选择。让更多企业和科研团队以可控成本用好智能体，让有限的计算资源支撑更多创新，为我国基础软件自主发展和产业智能化升级贡献力量。

\subsection{2.7 发展规划}\label{ux53d1ux5c55ux89c4ux5212}

PI OS 自 2026 年启动以来，围绕华为``高密部署智能体沙箱引擎''命题，已完成代码开发并形成阶段测试材料。项目将按照``研发验证---创业落地---市场拓展---持续发展''的路径，逐步由技术原型成长为面向企业交付的产品，推动科研成果与实际应用需求相结合。

\begin{longtable}[]{@{}
  >{\centering\arraybackslash}p{(\linewidth - 4\tabcolsep) * \real{0.1189}}
  >{\centering\arraybackslash}p{(\linewidth - 4\tabcolsep) * \real{0.1666}}
  >{\raggedright\arraybackslash}p{(\linewidth - 4\tabcolsep) * \real{0.7144}}@{}}
\caption*{表 2-1 PI OS 发展规划}\tabularnewline
\toprule\noalign{}
\rowcolor{WordTableHeader}
\begin{minipage}[b]{\linewidth}\centering
\TableHeaderText{发展阶段}
\end{minipage} & \begin{minipage}[b]{\linewidth}\centering
\TableHeaderText{时间}
\end{minipage} & \begin{minipage}[b]{\linewidth}\centering
\TableHeaderText{主要内容}
\end{minipage} \\
\midrule\noalign{}
\endfirsthead
\toprule\noalign{}
\rowcolor{WordTableHeader}
\begin{minipage}[b]{\linewidth}\centering
\TableHeaderText{发展阶段}
\end{minipage} & \begin{minipage}[b]{\linewidth}\centering
\TableHeaderText{时间}
\end{minipage} & \begin{minipage}[b]{\linewidth}\centering
\TableHeaderText{主要内容}
\end{minipage} \\
\midrule\noalign{}
\endhead
\bottomrule\noalign{}
\endlastfoot
\rowcolors{1}{WordTableStripe}{white}
准备期 & 2026年 & 完成项目立项、需求分析和代码开发，形成产品原型与阶段测试材料。 围绕企业指定环境开展适配，验证有限内存条件下的任务承载能力与运行稳定性。 完善部署工具、测试报告和使用文档，完成竞赛展示与创业方案论证。 \\
创业期 & 2027年 & 明确成果商业化授权范围，推进公司设立和创业团队建设。 面向企业研发团队、高校科研团队开展客户调研与小规模试点，根据反馈完善产品。 建立报价、交付和售后服务流程，探索部署适配、技术支持与年度维护等收入来源，争取形成首批付费客户案例。 \\
成长期 & 2028年至2030年 & 将试点经验转化为可复制的产品和交付方案，逐步拓展企业智能体平台及云服务商客户。 完善产品功能、稳定性和维护能力，支持更多实际任务场景。 建设市场与服务渠道，通过持续服务和客户续约，推动业务由单次项目交付向持续经营发展。 \\
稳定发展期 & 2031年及以后 & 持续跟进智能体应用需求，提升产品效率、可靠性和易用性，完善长期支持体系。 深化与国产操作系统、云服务及智能体应用生态的合作，扩大产品覆盖范围。 在国内形成稳定业务基础后，逐步探索海外合作与市场机会，向具有国际竞争力的基础软件品牌发展。 \\
\end{longtable}

